\documentclass[11pt,a4paper]{article}

\usepackage[utf8]{inputenc}
\usepackage[T1]{fontenc}
\usepackage[italian]{babel}
\usepackage[a4paper,margin=2.2cm]{geometry}
\usepackage{microtype}
\usepackage{xcolor}
\usepackage{hyperref}
\usepackage{enumitem}
\usepackage{listings}
\usepackage{tikz}
\usetikzlibrary{arrows.meta,positioning,shapes.geometric}

\definecolor{projectblue}{HTML}{1F4E79}
\definecolor{projectgreen}{HTML}{3A7D44}
\definecolor{projectgray}{HTML}{F3F5F7}
\definecolor{codegray}{HTML}{F7F7F7}

\hypersetup{
  colorlinks=true,
  linkcolor=projectblue,
  urlcolor=projectblue
}

\lstdefinestyle{terminal}{
  basicstyle=\ttfamily\small,
  backgroundcolor=\color{codegray},
  frame=single,
  framerule=0.2pt,
  breaklines=true,
  columns=fullflexible
}

\setlist[itemize]{leftmargin=1.2em,itemsep=0.25em}
\setlist[enumerate]{leftmargin=1.4em,itemsep=0.25em}

\title{\textbf{Project eBPF Runtime Monitor}\\
\large Stato MVP, demo operativa e prossimi passi}
\author{Simone Romano}
\date{Maggio 2026}

\begin{document}

\maketitle

\section{Obiettivo}

Il progetto realizza un runtime security monitor basato su eBPF, ispirato a
Tracee ma ridotto a un MVP piu' piccolo e controllabile. L'obiettivo attuale e'
dimostrare una pipeline end-to-end funzionante:

\begin{center}
\textit{build/embed oggetto eBPF $\rightarrow$ inizializzazione config
$\rightarrow$ selezione probe $\rightarrow$ load/attach $\rightarrow$ ring/perf buffer
$\rightarrow$ decoder $\rightarrow$ event guard $\rightarrow$ output leggibile}
\end{center}

Il target principale di sviluppo e test e' Rocky Linux 8.10 con kernel
\texttt{4.18.0-553.109.1.el8\_10.x86\_64}. Dove Tracee introduce componenti
production-grade piu' estesi, il progetto mantiene una struttura essenziale per
rendere chiaro il funzionamento dei blocchi fondamentali.

\section{Architettura attuale}

\begin{figure}[h]
\centering
\begin{tikzpicture}[
  node distance=0.42cm,
  box/.style={
    rectangle,
    rounded corners=2pt,
    draw=projectblue,
    fill=projectgray,
    text width=8.8cm,
    align=center,
    minimum height=0.78cm,
    font=\small
  },
  arrow/.style={-{Latex[length=2.2mm]}, thick, draw=projectblue}
]

\node[box] (build) {
  \textbf{Build artifact}\\
  \texttt{make bpf} genera \texttt{dist/project.bpf.o}; \texttt{go:embed} lo
  include nel binario Go
};
\node[box, below=of build] (init) {
  \textbf{CLI + initialize}\\
  \texttt{BPFObject()}: override da env/CLI oppure oggetto embedded; risolve
  anche il path BTF
};
\node[box, below=of init] (selection) {
  \textbf{Probe selection}\\
  \texttt{probes.Select(--events, --drop-events)} produce
  \texttt{selectedProbes} ed \texttt{eventSet}
};
\node[box, below=of selection] (runtime) {
  \textbf{Load + attach: \texttt{pkg/ebpf}}\\
  carica collection, apre ring buffer e perf buffer, poi attacca solo i probe selezionati
};
\node[box, below=of runtime] (kernel) {
  \textbf{Kernel data path}\\
  gli hook eBPF selezionati emettono eventi in \texttt{events\_ringbuf} o \texttt{events}
};
\node[box, below=of kernel] (decoder) {
  \textbf{Read + decode}\\
  \texttt{InitRingBuf}/\texttt{InitPerfBuf} $\rightarrow$ \texttt{bufferdecoder.DecodeEvent()}:
  raw bytes in contesto, nome evento e argomenti
};
\node[box, below=of decoder] (guard) {
  \textbf{Event guard userspace}\\
  \texttt{eventEnabled(event.EventName)} scarta eventuali eventi non abilitati
};
\node[box, below=of guard] (output) {
  \textbf{Output layer: \texttt{pkg/output}}\\
  \texttt{Printer}: formati \texttt{json} e \texttt{table}, JSON normalizzato,
  capability labels
};

\draw[arrow] (build) -- (init);
\draw[arrow] (init) -- (selection);
\draw[arrow] (selection) -- (runtime);
\draw[arrow] (runtime) -- (kernel);
\draw[arrow] (kernel) -- (decoder);
\draw[arrow] (decoder) -- (guard);
\draw[arrow] (guard) -- (output);

\end{tikzpicture}
\caption{Pipeline di inizializzazione e data path del tool.}
\end{figure}

La selezione eventi avviene in due momenti diversi. Prima dell'attach,
\texttt{probes.Select} decide quali programmi eBPF collegare al kernel. Dopo la
decodifica, \texttt{eventEnabled} resta come guardia userspace per evitare di
stampare eventi fuori dalla selezione configurata.

\section{Funzionalita' implementate}

\subsection{Loader e runtime eBPF}

Il runtime Go carica l'oggetto eBPF, apre sia la ring buffer sia il perf buffer
e mantiene i link agli hook kernel. Dopo il merge con il branch networking,
l'implementazione usa \texttt{github.com/aquasecurity/libbpfgo}, quindi il
Makefile prepara anche \texttt{libbpf} e i flag CGO necessari. I link vengono
chiusi in modo ordinato durante il cleanup.

\subsection{Probe registry e selezione eventi}

La lista degli hook supportati e' stata spostata in
\texttt{pkg/ebpf/probes/probes.go}. Ogni entry collega:

\begin{itemize}
  \item nome evento decodificato;
  \item nome del programma eBPF nell'oggetto compilato;
  \item hook kernel;
  \item tipo di attach, per esempio raw tracepoint o kprobe.
\end{itemize}

La CLI supporta:

\begin{lstlisting}[style=terminal]
--events sched_process_exec
--drop-events cap_capable
\end{lstlisting}

Questo permette di ridurre il rumore generato da hook molto frequenti come
\texttt{cap\_capable}.

\subsection{Decoder userspace}

Il package \texttt{pkg/bufferdecoder} trasforma i record binari prodotti dagli
hook eBPF in eventi Go. Il decoder conosce il protocollo condiviso con il lato C:
\texttt{event\_context\_t}, numero di argomenti e slot tipizzati per scalari e
stringhe.

\subsection{Output layer}

Il runtime non serializza piu' direttamente gli eventi. Invece usa
l'interfaccia \texttt{Printer} in \texttt{pkg/output}. Sono presenti due formati:

\begin{itemize}
  \item \texttt{json}: output machine-readable, una riga JSON per evento;
  \item \texttt{table}: output compatto per debug manuale da terminale.
\end{itemize}

Il JSON non espone piu' direttamente campi raw come array di byte. Campi come
\texttt{comm} e \texttt{uts\_name} vengono convertiti in stringhe. Gli argomenti
di tipo capability vengono arricchiti con label simboliche, ad esempio:

\begin{lstlisting}[style=terminal]
"value": 21, "label": "CAP_SYS_ADMIN"
\end{lstlisting}

\subsection{Oggetto eBPF embedded}

Seguendo il pattern di Tracee, l'oggetto \texttt{project.bpf.o} viene generato
in \texttt{dist/project.bpf.o} e incluso nel binario Go tramite
\texttt{go:embed}. Il path esplicito resta disponibile per debug, ma il comando
standard non richiede piu' \texttt{--bpf-object}.

\section{Demo operativa}

\subsection{Build}

\begin{lstlisting}[style=terminal]
cd demo_project
make build
\end{lstlisting}

Il target compila prima l'oggetto eBPF in \texttt{dist/project.bpf.o} e poi
costruisce il binario Go con l'oggetto embedded.

\subsection{Esecuzione con output compatto}

\begin{lstlisting}[style=terminal]
sudo ./dist/project --events sched_process_exec --output table
\end{lstlisting}

In un secondo terminale:

\begin{lstlisting}[style=terminal]
ls
whoami
bash -c 'echo demo'
\end{lstlisting}

Output atteso:

\begin{lstlisting}[style=terminal]
event=sched_process_exec pid=1234 tid=1234 uid=1000 comm=bash args=filename=/usr/bin/bash
\end{lstlisting}

\subsection{Esecuzione JSON normalizzata}

\begin{lstlisting}[style=terminal]
sudo ./dist/project --events cap_capable --output json
\end{lstlisting}

Esempio di output:

\begin{lstlisting}[style=terminal]
{"timestamp":4074609472044753,"event_name":"cap_capable","process":{"pid":1374462,"tid":1374462,"ppid":1348641,"uid":1000,"comm":"cpuUsage.sh","uts_name":"security-thesis"},"host":{"pid":1374462,"tid":1374462,"ppid":1348641},"kernel":{"syscall":56,"processor_id":0,"mnt_id":4026531840,"pid_id":4026531836},"args":[{"name":"cap","type":1,"value":21,"label":"CAP_SYS_ADMIN"}]}
\end{lstlisting}

Per ridurre il rumore:

\begin{lstlisting}[style=terminal]
sudo ./dist/project --drop-events cap_capable --output table
\end{lstlisting}

\section{Differenze rispetto a Tracee}

Tracee e' un progetto production-grade con policy engine, filtering lato kernel,
pipeline eventi completa, enrichment, detection e supporto kernel piu' ampio.
Questo progetto riprende alcuni pattern, ma li mantiene in forma MVP:

\begin{itemize}
  \item registry statico dei probe invece di un sistema completo di
  \texttt{ProbeGroup};
  \item selezione eventi iniziale invece di policy engine completo;
  \item decoder custom per il protocollo ridotto del progetto;
  \item runtime basato su \texttt{libbpfgo}, ma senza ancora tutta
  l'infrastruttura di Tracee per policy, enrichment e detection;
  \item output layer semplificato, ispirato al ruolo del sink stage;
  \item oggetto eBPF embedded nel binario, come in Tracee.
\end{itemize}

\section{Prossimi passi}

Dalla call del 7 maggio e dallo stato tecnico attuale emergono questi step:

\begin{enumerate}
  \item tornare sulla parte eBPF e completare gli hook;
  \item studiare il filtering lato kernel usato da Tracee, in particolare policy
  mask e mappe eBPF di configurazione;
  \item analizzare gli hook asincroni e la correlazione tra evento e processo;
  \item arricchire l'output con mapping di syscall, resource limit e opzioni
  \texttt{prctl};
  \item introdurre una prima detection logic sopra la pipeline eventi.
\end{enumerate}

\section{Comandi di verifica}

\begin{lstlisting}[style=terminal]
make build
make help
make run_table
make filtered
\end{lstlisting}

\end{document}
